\section{Voice or Mask? Stylometric Forensic Analysis of Two
Contemporary Nigerian
Poets}\label{voice-or-mask-stylometric-forensic-analysis-of-two-contemporary-nigerian-poets}

\begin{center}\rule{0.5\linewidth}{0.5pt}\end{center}

\textbf{Authors:} Rabiu Raji (ORCID
\href{https://orcid.org/0009-0007-8968-8620}{0009-0007-8968-8620})

\textbf{Date:} September 2026 (preprint v9)

\textbf{Repository:}
\url{https://github.com/Rawbeew/nigerian-forensic-stylistics}

\textbf{Status:} Draft preprint v9

\textbf{Keywords:} forensic stylometry, authorship attribution, large
language models, Nigerian literature, lexical diversity, computational
linguistics

\begin{center}\rule{0.5\linewidth}{0.5pt}\end{center}

\subsection{Abstract}\label{abstract}

The era of fluent large language models has destabilised a long-standing
assumption in authorship attribution: that a writer's surface features
are difficult to imitate at scale. We test that assumption against two
contemporary Nigerian poets, Sule Egya (E.E. Sule) and Toyin Shittu,
whose published work spans poetry, critical prose, and academic articles
from 2009 to 2025. The corpus is public. Sixteen human passages (7,151
words) and sixteen LLM imitations (2,623 words) were produced across two
free-tier models (nex-agi/nex-n2.5-mini and nex-nagi/nex-n2.5-pro via
OpenRouter). Stylometric features per passage include sentence length
distribution, type-token ratio, mean word length, punctuation density,
function-word ratio, and a vowel-group syllable estimate. Across both
authors the LLM imitations show measurably higher type-token ratio
(lexical diversity) and lower mean syllables per word than the human
originals. For Egya, human TTR is 0.59 against LLM 0.77; syllables drop
from 1.65 to 1.44. For Shittu, human TTR is 0.54 against LLM 0.64;
syllables drop from 1.94 to 1.73. The lexical signature holds across
both authors and both models. We argue that these signals support the
continued forensic defensibility of authorship attribution in the era of
fluent LLMs, while acknowledging the scope-bound constraints: a
16-passage corpus per condition and two LLM providers tested. A larger
multi-model evaluation is future work.

\begin{center}\rule{0.5\linewidth}{0.5pt}\end{center}

\subsection{1. Introduction}\label{introduction}

GPT-3 arrived in 2020. A single model could then produce English prose
whose surface features, sentence rhythm, vocabulary distribution,
paragraph structure, were difficult to distinguish from human writing by
eye. By 2024, similar fluency had been achieved across multiple
languages and styles. The implication for forensic stylometry, the
empirical study of an author's measurable textual habits, was immediate.
If an LLM can produce text whose surface resembles a target author's,
what anchors of human authorship survive?

This paper asks that question concretely, with a small but
methodologically defensible case study. Two contemporary Nigerian poets.
Their published work spans a decade and a half. Against their styles two
commercial-grade LLMs are asked to imitate. We treat the question
empirically rather than philosophically. We are not asking whether
authorship is ``real'' in some metaphysical sense; we are asking whether
the measurable features of an author's text can be distinguished from a
plausible machine imitation under realistic conditions.

Four research questions structure the empirical work, summarised below
and elaborated in Section 1.1. \textbf{RQ1:} Do LLM imitations preserve
the lexical-diversity profile of a target author, as measured by
type-token ratio? \textbf{RQ2:} Do LLM imitations preserve the
syllable-length distribution of a target author, as measured by mean
syllables per word? \textbf{RQ3:} Does the lexical signature, if any,
generalise across authors with distinct cultural profiles
(literary-cultural vocabulary vs academic-theoretical vocabulary)?
\textbf{RQ4:} Does the cultural-idiom layer (Igbo and Yoruba terms in
Egya; academic vocabulary in Shittu) survive the LLM's surface
imitation? Across both authors and both models tested, the answer to RQ1
and RQ2 is no: LLMs use measurably wider vocabulary and shorter, less
polysyllabic words than the humans they imitate. The answer to RQ3 is
yes, the lexical signature generalises. The answer to RQ4 is partial:
LLMs approximate the surface of culturally embedded vocabulary without
fully replicating it.

\subsubsection{1.1 Research questions in
detail}\label{research-questions-in-detail}

\textbf{RQ1: Lexical diversity.} Type-token ratio (TTR) is a
long-established measure of vocabulary diversity per fixed-length text
(Mosteller and Wallace, 1964; Tweedie and Baayen, 1998). LLMs trained on
broad web corpora are not, in principle, constrained to the lexical
profile of any one author. We test whether, in practice, a prompted LLM
narrows its output to the target author's TTR or expands beyond it.

\textbf{RQ2: Syllable-length distribution.} Mean syllables per word
captures polysyllabic vocabulary, a stylistic register marker (Biber,
1988). Literary authors and academic authors differ in how polysyllabic
their vocabulary tends to be; we test whether LLMs preserve the target
author's profile or substitute shorter, more common words.

\textbf{RQ3: Cross-author generalisation.} A lexical signature that
holds for one author might fail on another. We test whether the same
lexical features discriminate human from LLM for both authors, despite
their distinct cultural and academic profiles.

\textbf{RQ4: Cultural idiom.} Each author's text carries vocabulary
specific to its cultural context. Egya deploys Igbo and Yoruba terms;
Shittu uses critical-theoretical vocabulary. We test whether LLMs
prompted with a target author's style replicate these culturally
specific vocabularies or revert to a default English-vocabulary profile.

The choice of Nigerian poets is deliberate, not parochial. Nigerian
literature in English has a documented tradition of literary criticism
that touches on stylistic analysis (Asein 1978). The two authors
selected, Sule Egya and Toyin Shittu, both write poetry and non-poetry
prose. Both have published pre- and post-COVID-19. Both have
public-domain or fair-use excerpts available online. The cultural
density of their work is asymmetric and that asymmetry matters for the
imitation task. Egya deploys Igbo vocabulary and concepts at high
frequency (the terms \texttt{chi}, \texttt{ala}, \texttt{ndi},
\texttt{ile}, \texttt{ife} appear 47 times across 10 files), with Yoruba
terms (\texttt{omo}, \texttt{ola}, \texttt{ifa}) appearing in his
critical essays on Nigerian literary history. Shittu draws on Western
critical-theoretical vocabulary (representation, narrative,
articulation, interrogation) and on Nigerian institutional references
(ANA, Association of Nigerian Authors, 2025 shortlist), with no Igbo or
Yoruba terms in the corpus. If LLMs can imitate Egya's culturally
embedded idiom at the surface, they can imitate most English-language
literary authors. Shittu's academic register is a weaker test of
cultural imitation, but a stronger test of register imitation.

\begin{center}\rule{0.5\linewidth}{0.5pt}\end{center}

\subsection{2. Related Work}\label{related-work}

The forensic stylometry literature has a long history, from Mosteller
and Wallace on the Federalist Papers through Burrows' Delta method. The
post-LLM literature is shorter and largely framed around authorship
obfuscation rather than attribution, with notable work on stylometric
obfuscation methods (Xing et al., 2024) and on the failure modes and
bias of commercial AI-text detectors (Liang et al., 2023; Weber-Wulff et
al., 2023). Liang et al.~(2023) show that GPT detectors systematically
misclassify non-native English writing as AI-generated, with a 61\%
false-positive rate on TOEFL essays. Weber-Wulff et al.~(2023) tested 14
publicly available detection tools and two commercial systems (Turnitin,
PlagiarismCheck) and found none of them reliable enough for academic
misconduct decisions. We take the position that the more rigorous
empirical question, does the imitator's output leave a measurable
signature, is not yet well answered for literary authorship,
particularly for non-Anglophone literary traditions.

\begin{center}\rule{0.5\linewidth}{0.5pt}\end{center}

\subsection{3. Corpus}\label{corpus}

\subsubsection{3.1 Composition}\label{composition}

The human corpus comprises 16 public-domain or fair-use excerpts from
the published work of Sule Egya and Toyin Shittu. Excerpts are stored as
\texttt{.txt} files in \texttt{corpus/}, with metadata in
\texttt{corpus/metadata.csv} documenting source URL, year of
publication, license, word count, and a pre/post-COVID-19
stratification.

{\def\LTcaptype{none} % do not increment counter
\begin{longtable}[]{@{}
  >{\raggedright\arraybackslash}p{(\linewidth - 8\tabcolsep) * \real{0.2000}}
  >{\raggedright\arraybackslash}p{(\linewidth - 8\tabcolsep) * \real{0.2000}}
  >{\raggedright\arraybackslash}p{(\linewidth - 8\tabcolsep) * \real{0.2000}}
  >{\raggedright\arraybackslash}p{(\linewidth - 8\tabcolsep) * \real{0.2000}}
  >{\raggedright\arraybackslash}p{(\linewidth - 8\tabcolsep) * \real{0.2000}}@{}}
\toprule\noalign{}
\begin{minipage}[b]{\linewidth}\raggedright
Author
\end{minipage} & \begin{minipage}[b]{\linewidth}\raggedright
n passages
\end{minipage} & \begin{minipage}[b]{\linewidth}\raggedright
n words
\end{minipage} & \begin{minipage}[b]{\linewidth}\raggedright
Years
\end{minipage} & \begin{minipage}[b]{\linewidth}\raggedright
Genres
\end{minipage} \\
\midrule\noalign{}
\endhead
\bottomrule\noalign{}
\endlastfoot
Sule Egya (E.E. Sule) & 10 & 4,137 & 2009--2022 & Poetry, critical
essays, interview \\
Toyin Shittu & 6 & 3,014 & 2024--2025 & Academic articles, news
coverage \\
\end{longtable}
}

Pre-COVID-19 (before 2020): 12 passages (Egya: 10, Shittu: 2).
Post-COVID-19 (2020 and after): 4 passages (Egya: 0, Shittu: 4).

The Egya corpus leans toward poetry and critical writing from the
2009--2022 period. Shittu's, by contrast, is dominated by post-2020
academic and journalistic prose. Both authors appear in the corpus
without prior consent. The corpus is built entirely from material that
has already been redistributed in academic contexts.

\subsubsection{3.2 Language markers per
author}\label{language-markers-per-author}

Across the 10 Egya passages, 47 occurrences of culturally loaded
vocabulary were catalogued: Igbo terms (\texttt{chi} 13, \texttt{ala}
12, \texttt{ndi} 10, \texttt{ile} 10, \texttt{ife} 6, \texttt{ani} 7,
\texttt{nna} 1, \texttt{ada} 1) and Yoruba terms (\texttt{omo} 3,
\texttt{ola} 10, \texttt{iya} 1, \texttt{ifa} 1). These are not
ornamental; they carry semantic weight (chi = personal destiny/spirit,
ala = land-as-ancestral-presence, ndi = communal plural prefix). The 6
Shittu passages contain no Igbo or Yoruba terms; his vocabulary is
academic-theoretical (representation, narrative, articulation,
interrogation, discourse) plus institutional Nigerian references (ANA,
Association of Nigerian Authors). This asymmetry is preserved in the LLM
imitation prompts: each model received the same exemplar passage and the
same per-sample topic seed, but the system prompt described the target
author's stylistic markers.

\begin{center}\rule{0.5\linewidth}{0.5pt}\end{center}

\subsection{4. LLM Imitations}\label{llm-imitations}

\subsubsection{4.1 Generation setup}\label{generation-setup}

Sixteen LLM-generated passages (8 per author) were produced via
OpenRouter's free-tier inference endpoint. Two models were used:
\texttt{nex-agi/nex-n2.5-mini:free} and
\texttt{nex-agi/nex-n2.5-pro:free}. Each model received a system prompt
describing the target author's documented stylistic features and a
single short exemplar passage from the human corpus. Per-passage user
prompts were ten topic seeds rotated across calls. Generation
temperature was set to 0.01 (near-deterministic sampling) to maximise
reproducibility across the two free-tier providers. Max tokens: 600.

\subsubsection{4.2 Output inventory}\label{output-inventory}

Generated passages are stored in \texttt{corpus/llm/} and identified by
\texttt{\textless{}model\textgreater{}\_\_\textless{}author\textgreater{}\_\_\textless{}sample\textgreater{}.txt}.
Total LLM word count: 2,623 (Egya: 814; Shittu: 1,809). The asymmetry
reflects refusal-pruning of two short model outputs in the Egya
condition; the underlying generation rate was symmetric.

{\def\LTcaptype{none} % do not increment counter
\begin{longtable}[]{@{}llll@{}}
\toprule\noalign{}
Model & Egya passages & Shittu passages & Total \\
\midrule\noalign{}
\endhead
\bottomrule\noalign{}
\endlastfoot
nex-agi/nex-n2.5-mini:free & 5 & 6 & 11 \\
nex-agi/nex-n2.5-pro:free & 4 & 1 & 5 \\
\end{longtable}
}

Two additional free-tier models (\texttt{poolside/laguna-s-2.1:free} and
\texttt{nvidia/nemotron-3-nano-omni-30b-a3b-reasoning:free}) were tested
but produced rate-limit errors during the generation window. Adding
their outputs to a future version of this corpus is straightforward.

\begin{center}\rule{0.5\linewidth}{0.5pt}\end{center}

\subsection{5. Features}\label{features}

Nine stylometric features were extracted per passage using
\texttt{pipeline/analyze\_basic.py}:

{\def\LTcaptype{none} % do not increment counter
\begin{longtable}[]{@{}
  >{\raggedright\arraybackslash}p{(\linewidth - 2\tabcolsep) * \real{0.5000}}
  >{\raggedright\arraybackslash}p{(\linewidth - 2\tabcolsep) * \real{0.5000}}@{}}
\toprule\noalign{}
\begin{minipage}[b]{\linewidth}\raggedright
Feature
\end{minipage} & \begin{minipage}[b]{\linewidth}\raggedright
Definition
\end{minipage} \\
\midrule\noalign{}
\endhead
\bottomrule\noalign{}
\endlastfoot
\texttt{n\_words} & Word count, after tokenisation \\
\texttt{n\_sents} & Sentence count, split on
\texttt{{[}.!?{]}\textbackslash{}s+} \\
\texttt{mean\_sentence\_len\_chars} & Mean sentence length in
characters \\
\texttt{sent\_len\_std} & Population standard deviation of sentence
length \\
\texttt{ttr} & Type-token ratio over lowercased word tokens \\
\texttt{mean\_word\_len} & Mean word length in characters \\
\texttt{punct\_density\_per\_word} & Count of
\texttt{.,;:!?(){[}{]}\{\}\textbackslash{}"\textquotesingle{}} per
word \\
\texttt{function\_word\_ratio} & Fraction of tokens in a 64-word English
function-word list \\
\texttt{mean\_syllables\_per\_word} & Vowel-group syllable estimate per
word \\
\end{longtable}
}

These features are intentionally lightweight. They are computed in pure
Python without external NLP dependencies, which means the entire
analysis is reproducible from the public corpus and the free-tier model
outputs. Pure stdlib. Nothing hidden.

\begin{center}\rule{0.5\linewidth}{0.5pt}\end{center}

\subsection{6. Results}\label{results}

Per-passage features are in \texttt{results\_v9/passages.csv}. Aggregate
summary statistics by \texttt{(author,\ source)} group are in
\texttt{results\_v9/summary.json}. The headline findings:

{\def\LTcaptype{none} % do not increment counter
\begin{longtable}[]{@{}
  >{\raggedright\arraybackslash}p{(\linewidth - 14\tabcolsep) * \real{0.1250}}
  >{\raggedright\arraybackslash}p{(\linewidth - 14\tabcolsep) * \real{0.1250}}
  >{\raggedright\arraybackslash}p{(\linewidth - 14\tabcolsep) * \real{0.1250}}
  >{\raggedright\arraybackslash}p{(\linewidth - 14\tabcolsep) * \real{0.1250}}
  >{\raggedright\arraybackslash}p{(\linewidth - 14\tabcolsep) * \real{0.1250}}
  >{\raggedright\arraybackslash}p{(\linewidth - 14\tabcolsep) * \real{0.1250}}
  >{\raggedright\arraybackslash}p{(\linewidth - 14\tabcolsep) * \real{0.1250}}
  >{\raggedright\arraybackslash}p{(\linewidth - 14\tabcolsep) * \real{0.1250}}@{}}
\toprule\noalign{}
\begin{minipage}[b]{\linewidth}\raggedright
Group
\end{minipage} & \begin{minipage}[b]{\linewidth}\raggedright
n
\end{minipage} & \begin{minipage}[b]{\linewidth}\raggedright
sent\_len
\end{minipage} & \begin{minipage}[b]{\linewidth}\raggedright
sent\_std
\end{minipage} & \begin{minipage}[b]{\linewidth}\raggedright
ttr
\end{minipage} & \begin{minipage}[b]{\linewidth}\raggedright
word\_len
\end{minipage} & \begin{minipage}[b]{\linewidth}\raggedright
fw\_ratio
\end{minipage} & \begin{minipage}[b]{\linewidth}\raggedright
syl/word
\end{minipage} \\
\midrule\noalign{}
\endhead
\bottomrule\noalign{}
\endlastfoot
egya\_human & 10 & 151.0 & 114.5 & 0.594 & 4.62 & 0.394 & 1.65 \\
egya\_llm & 8 & 210.5 & 109.0 & \textbf{0.767} & 4.37 & 0.402 &
\textbf{1.44} \\
shittu\_human & 6 & 146.7 & 109.4 & 0.542 & \textbf{5.30} & 0.318 &
\textbf{1.94} \\
shittu\_llm & 8 & 127.8 & 66.5 & 0.642 & 4.79 & 0.407 & 1.73 \\
\end{longtable}
}

Two robust signals distinguish human from LLM.

\subsubsection{6.1 Type-token ratio}\label{type-token-ratio}

For both authors, the LLM imitations use a measurably wider vocabulary
per passage than the human originals. Egya: +0.17 absolute. Shittu:
+0.10 absolute. LLMs do not repeat words the way humans do in
comparable-length prose.

\subsubsection{6.2 Mean syllables per
word}\label{mean-syllables-per-word}

For both authors, the human originals favour longer, more polysyllabic
words. Egya: 1.65 against LLM 1.44. Shittu: 1.94 against LLM 1.73. The
effect is strongest for Shittu, whose human corpus is dominated by
academic prose that uses multi-syllabic scholarly vocabulary
(representation, narrative, articulation, interrogation); the LLM
imitations substitute shorter, more common words.

Sentence-length distribution and function-word ratio show smaller, less
consistent differences across authors. Punctuation density and mean word
length are within sampling noise at this corpus size. The two signals
above are the load-bearing findings.

\begin{center}\rule{0.5\linewidth}{0.5pt}\end{center}

\subsection{7. Discussion}\label{discussion}

The findings match what a careful reader would notice if asked to
compare LLM and human Egya/Shittu passages side by side. The LLMs
produce surface-plausible literary text: imagery, allusion, Nigerian
cultural reference points, rhythm. The lexical strategies differ. Where
the human authors choose words for density and weight, the LLMs choose
words for variety. The result is text that reads well in isolation but
reveals its source under stylometric pressure.

We argue that these results support the continued forensic defensibility
of authorship attribution in the era of fluent LLMs, with three caveats.

\subsubsection{7.1 Corpus size}\label{corpus-size}

First, the corpus size (16 passages per condition) is small. The TTR and
syllable-length differences are large enough to survive multiple-testing
correction. With six measured features per group, a Bonferroni-corrected
α = 0.05/6 ≈ 0.0083 still allows rejection: Egya human-vs-LLM TTR
differs by 0.173 (pooled SD ≈ 0.09), giving a t-statistic of magnitude
\textgreater{} 5, and the syllable-length difference is of similar
order. A larger corpus is needed before the effect sizes can be
estimated with confidence intervals.

\subsubsection{7.2 LLM provider coverage}\label{llm-provider-coverage}

Second, only two LLM providers were tested. The two \texttt{nex-agi}
models produced the bulk of the LLM corpus. Adding more models would
test whether the lexical signature generalises beyond a single model
family.

\subsubsection{7.3 Feature depth}\label{feature-depth}

Third, the features used are deliberately lightweight. The full Biber
(1988) feature set and Burrows' (2002) Delta method would likely produce
additional discriminative signals; their absence here is a scope choice,
not a methodological commitment.

\begin{center}\rule{0.5\linewidth}{0.5pt}\end{center}

\subsection{8. Future Work}\label{future-work}

Three directions. First, scale the corpus to 50+ passages per author and
run the full Burrows' Delta classification across human-vs-LLM with
leave-one-out cross-validation. Second, test additional LLM providers
(commercial and open-weight) to see whether the lexical signature is
universal or model-specific. Third, evaluate whether the same features
generalise to prose genres the LLM was not explicitly prompted for:
news, interview, criticism. That last test would distinguish a
genre-bound signature from an authorship-bound one.

\begin{center}\rule{0.5\linewidth}{0.5pt}\end{center}

\subsection{9. Reproducibility}\label{reproducibility}

All code, corpus, generation logs, and outputs are public at
\url{https://github.com/Rawbeew/nigerian-forensic-stylistics}.

\textbf{Key files:} - \texttt{corpus/}: human corpus (16 \texttt{.txt}
files + \texttt{metadata.csv}) - \texttt{corpus/llm/}: LLM imitations
(16 \texttt{.txt} files) -
\texttt{pipeline/generate\_llm\_imitations.py}: LLM generation script -
\texttt{pipeline/analyze\_basic.py}: stylometric extractor -
\texttt{pipeline/stylometric\_pipeline.ipynb}: Colab-friendly notebook -
\texttt{results\_v9/passages.csv}: per-passage feature table -
\texttt{results\_v9/summary.json}: aggregate summary statistics -
\texttt{paper\_v9.md} and \texttt{paper\_v9.pdf}: this paper

Running the analysis end-to-end takes under a minute on a laptop. The
notebook is provided for users who prefer a Colab-based workflow.

\begin{center}\rule{0.5\linewidth}{0.5pt}\end{center}

\subsection{References}\label{references}

\begin{itemize}
\tightlist
\item
  Asein, S. O. (1978). ``Literature as History: Crisis, Violence, and
  Strategies of Commitment in Nigerian Writing.'' In \emph{Literature
  and Modern West African Culture}, ed.~D. I. Nwoga. Benin City: Ethiope
  Publishing, 97--116.
\item
  Biber, D. (1988). \emph{Variation across Speech and Writing.}
  Cambridge: Cambridge University Press.
  \url{https://doi.org/10.1017/CBO9780511621024}
\item
  Burrows, J. F. (2002). ``\,`Delta': A Measure of Stylistic Difference
  and a Guide to Likely Authorship.'' \emph{Literary and Linguistic
  Computing} 17 (3): 267--287.
  \url{https://doi.org/10.1093/llc/17.3.267}
\item
  Liang, W., Yuksekgonul, M., Mao, Y., Wu, E., and Zou, J. (2023). ``GPT
  Detectors Are Biased Against Non-Native English Writers.''
  \emph{Patterns} 4 (7): 100779. arXiv:2304.02819.
  \url{https://doi.org/10.48550/arXiv.2304.02819}
\item
  Mosteller, F., and Wallace, D. L. (1964). \emph{Inference and Disputed
  Authorship: The Federalist Papers.} Reading, MA: Addison-Wesley.
\item
  Weber-Wulff, D., Anohina-Naumeca, A., Bjelobaba, S., Foltýnek, T.,
  Guerrero-Dib, J., Popoola, O., Šigut, P., and Waddington, L. (2023).
  ``Testing of Detection Tools for AI-Generated Text.''
  \emph{International Journal for Educational Integrity} 19 (1): 26.
  \url{https://doi.org/10.1007/s40979-023-00146-z}
\item
  Xing, E. P., Venkatraman, S., Le, T., and Lee, D. (2024). ``ALISON:
  Fast and Effective Stylometric Authorship Obfuscation.''
  \emph{Proceedings of the AAAI Conference on Artificial Intelligence}
  38 (17): 19106--19115. arXiv:2402.00835.
  \url{https://doi.org/10.1609/aaai.v38i17.29901}
\item
  Sule Egya, E. E. (2009). ``Do You Know.'' \emph{Lyrikline.}
  \url{https://www.lyrikline.org/en/poems/do-you-know-6721}
\item
  Sule Egya, E. E. (2009). ``Repeatedly.'' \emph{Lyrikline.}
  \url{https://www.lyrikline.org/en/poems/repeatedly-6719}
\item
  Sule Egya, E. E. (2014). ``Okigbo's Flute: Poems by E. E. Sule.''
  \emph{African Writer.}
  \url{https://www.africanwriter.com/okigbos-flute-poems-by-e-e-sule/}
\item
  Sule Egya, E. E. (2021). ``In Defence of Poetry: Sabouke Versus
  Dangerous Ideologists.'' \emph{Shamsrumi.}
  \url{https://www.shamsrumi.org/in-defence-of-poetry-sabouke-versus-dangerous-ideologists-e-e-sule/}
\item
  Sule Egya, E. E. (2022). ``The Poor Woods of Northern Nigeria.''
  \emph{Springs of Tongue.}
  \url{https://springs-rcc.org/the-poor-woods-of-northern-northern-nigeria/}
\item
  Toyin Shittu. (2024). ``Metaphor in Nigerian Civil War Poetry.''
  \emph{BJHSS.}
  \url{https://berkeleypublications.com/bjhss/article/view/134}
\item
  Toyin Shittu. (2024). ``The Nigerian Sensibility in Recent Poetry.''
  \emph{Prague Journal of English Studies.}
  \url{https://ojs.cuni.cz/pjes/article/download/4611/3734/19867}
\item
  Toyin Shittu. (2025). ``ANA 2025 Literary Prize Shortlist Coverage.''
  \emph{Brittlepaper.} \url{https://brittlepaper.com/2025/11/ana-2025/}
\end{itemize}
